\chapter{Introduction}
\label{ch:introduction}

University students constantly need to ask questions, share study materials, and
coordinate around exams and deadlines. In practice, most of this communication is
scattered across social media groups and instant messaging channels that were never
designed for academic collaboration. Answers to the same question are posted again
and again, valuable explanations sink to the bottom of a chat history within hours,
course materials are shared as loose attachments that are hard to find later, and
there is no way to tell a reliable answer from a guess. As a result, knowledge that
one cohort of students has already worked out is rarely available to the next one.

\emph{TalkCS} is a web-based academic forum platform that addresses this problem. It
gives a university a single, structured place where students can ask and answer
questions, upload and rate course resources, follow each other, chat in real time,
and keep track of important academic dates. Discussions are organized into
categories that correspond to courses or subjects, every post can be tagged and
searched, and a voting and reputation system rewards helpful contributions so that
the best answers rise to the top. The platform is designed to be deployed by an
institution and customized to its needs: administrators manage categories, user
roles and branding, and can monitor the activity of the community through a simple
analytics dashboard.

The system is built as a classic three-tier web application. The backend is a
REST and WebSocket API implemented in Java with the Spring Boot framework~\cite{springboot},
secured with Spring Security~\cite{springsecurity} and stateless JSON Web Tokens~\cite{jwt},
and stores its data in a PostgreSQL~\cite{postgresql} relational database. The frontend is a
single-page application written in TypeScript~\cite{typescript} with React~\cite{react}, built with
Vite~\cite{vite} and styled with Tailwind CSS~\cite{tailwind}. Real-time features (live comment
threads and direct messaging) use the STOMP protocol over WebSocket connections~\cite{stomp}.

The main features of TalkCS are:

\begin{itemize}
  \item \textbf{Question-and-answer discussions:} Users publish posts with Markdown
    formatting and syntax-highlighted code blocks into subject categories, attach up to
    five tags, and receive automatic tag suggestions based on the text they write.
  \item \textbf{Duplicate question detection:} While a user is typing a new post, the
    system searches the selected category for similar existing posts using a ranking
    algorithm that combines the BM25 relevance function~\cite{bm25} with title and tag
    similarity, and warns the user before a duplicate is created.
  \item \textbf{Threaded comments and accepted answers:} Posts can be answered with
    arbitrarily nested comments. The author of a question can accept one answer, which
    marks the question as \emph{solved}.
  \item \textbf{Voting, reputation and gamification:} Posts, comments and resources can
    be up- or down-voted. Votes and accepted answers change the reputation of the
    author, which determines the user's level, the leaderboard position, and a set of
    milestone and subject-expertise badges.
  \item \textbf{Resource library:} Each category has a dedicated library where users
    upload, download and vote on course materials independently of the discussions.
  \item \textbf{Social features:} User profiles, a follow system, \code{@username}
    mentions, bookmarks, in-app notifications, and a daily e-mail digest.
  \item \textbf{Real-time direct messaging} between any two users of the platform.
  \item \textbf{Shared academic calendar:} Users keep private events and propose public
    ones (exams, deadlines, lectures), which professors or administrators review.
  \item \textbf{Administration:} Usage analytics, user and role management, category
    management with archiving, and branding configuration.
  \item \textbf{Comfortable user experience:} Global search, trending posts and
    resources, a responsive layout for mobile devices, and dark and light themes.
\end{itemize}

The rest of this thesis is organized as follows. Chapter~\ref{ch:userdoc} is the user
documentation: it describes the system requirements, the installation process, and
walks through every feature of the application from the point of view of its users.
Chapter~\ref{ch:devdoc} is the developer documentation, covering the architecture,
the database design, the structure of the backend and frontend code, and the most
interesting implementation details. Chapter~\ref{ch:testing} presents the testing
strategy and its results, and Chapter~\ref{ch:conclusion} summarizes the work and
outlines possible future improvements.
